\chapter{FUNDAMENTAÇÃO TEÓRICA}
\label{cap:fundamentacao}

\section{Aprendizagem Matemática e Evidências de Desempenho}

Aprender matemática não se limita a memorizar procedimentos ou obter respostas corretas. O processo envolve construir significados, relacionar novos conceitos ao conhecimento prévio, desenvolver formas de representação, justificar estratégias e transferir o que foi aprendido para problemas diferentes \cite{Piaget1972EN, Vygotsky1978, Ausubel1968}. Por essa razão, uma medida isolada de acerto, nota ou tempo de resposta oferece apenas uma evidência parcial do processo de aprendizagem.

A Base Nacional Comum Curricular (BNCC) associa o letramento matemático às capacidades de raciocinar, representar, comunicar e argumentar matematicamente, além de formular e resolver problemas em contextos variados \cite{BNCC2018}. Essa definição reforça que competência matemática envolve a mobilização integrada de conhecimentos, procedimentos e atitudes, e não apenas o desempenho observado em uma atividade.

Para evitar que medidas distintas fossem tratadas como equivalentes, este trabalho adotou definições operacionais. A avaliação educacional parte de observações produzidas por tarefas e respostas para realizar inferências sobre conhecimentos e habilidades; essas inferências dependem do modelo de cognição, dos instrumentos e dos procedimentos de mensuração empregados \cite{NRC2001}. Em avaliações como o PISA, a proficiência é estimada em uma escala e interpretada por níveis que descrevem o que os estudantes tipicamente conseguem realizar \cite{OECDPISA2022Framework}.

\textbf{Desempenho observado} corresponde ao registro diretamente obtido em uma tarefa, como acerto, nota, número de tentativas, estratégia apresentada ou tempo de resolução \cite{NRC2001}. \textbf{Proficiência estimada} corresponde a uma inferência sobre a posição do estudante em uma escala de conhecimentos e habilidades, calculada a partir de um conjunto de evidências e de um modelo de mensuração \cite{NRC2001, OECDPISA2022Framework}. \textbf{Competência} refere-se à capacidade de mobilizar conhecimentos, procedimentos, estratégias e atitudes para enfrentar situações e resolver problemas \cite{BNCC2018}. Por sua vez, \textbf{aprendizagem} é entendida como um processo de transformação construído ao longo do tempo, evidenciado por mudanças na compreensão, na autonomia, nas estratégias e na capacidade de aplicar conhecimentos em novos contextos \cite{Piaget1972EN, Vygotsky1978, Ausubel1968}.

Essas definições são relacionadas, mas não intercambiáveis. Um modelo computacional pode reconhecer padrões nos dados de desempenho e produzir estimativas de proficiência. Entretanto, não observa diretamente todos os processos cognitivos, sociais e afetivos envolvidos na aprendizagem. Suas saídas devem ser interpretadas como apoio à análise docente, e não como diagnóstico definitivo do estudante.

\section{Contribuições de Teorias de Aprendizagem}

\subsection{Construção Ativa do Conhecimento}

A epistemologia genética de Piaget descreve o conhecimento como uma construção realizada pelo sujeito em interação com o meio \cite{Piaget1972EN}. Situações que desafiam explicações anteriores podem provocar reorganizações cognitivas. No ensino de matemática, essa perspectiva valoriza a resolução de problemas, a comparação de estratégias e a compreensão dos conceitos, em contraste com a repetição mecânica de algoritmos.

\subsection{Mediação, Zona de Desenvolvimento Proximal e Andaimento Pedagógico}

Vygotsky definiu a Zona de Desenvolvimento Proximal como a distância entre aquilo que o estudante realiza de forma autônoma e aquilo que consegue realizar com mediação \cite{Vygotsky1978}. O conceito de andaimento pedagógico (\textbf{\textit{scaffolding}}) foi formulado posteriormente por Wood, Bruner e Ross para descrever um suporte temporário oferecido durante a resolução de problemas e retirado gradualmente conforme aumenta a autonomia do aprendiz \cite{WoodBrunerRoss1976}. Sistemas computacionais podem apoiar esse processo de mediação, mas a adequação pedagógica das intervenções depende do contexto e da interpretação do professor.

\subsection{Conhecimento Prévio e Aprendizagem Significativa}

Ausubel destacou que novas informações são aprendidas de forma significativa quando se relacionam a conhecimentos relevantes já existentes na estrutura cognitiva do estudante \cite{Ausubel1968}. Para a matemática, isso implica identificar concepções anteriores, dificuldades e relações entre conteúdos. Um sistema de apoio pode organizar evidências sobre esses aspectos, mas não deve confundir ausência de acerto com ausência de compreensão sem considerar a natureza da tarefa e a estratégia utilizada.

\subsection{Avaliação Formativa e Retorno Pedagógico}

A avaliação formativa utiliza evidências produzidas durante o processo de aprendizagem para orientar decisões sobre as próximas ações de ensino e estudo \cite{BlackWiliam1998}. A efetividade do retorno pedagógico depende de seu foco, de sua clareza e da forma como ajuda o estudante a compreender os objetivos, avaliar seu estado atual e identificar ações para avançar \cite{HattieTimperley2007}. Em contexto de aprendizagem matemática, Mertasari et al. observaram diferenças no desenvolvimento de habilidades metacognitivas entre modalidades de avaliação formativa, oferecendo uma evidência empírica específica para esse domínio \cite{Performance2023_014}.

Uma previsão algorítmica, por si só, não constitui retorno formativo. Para cumprir essa função, a saída precisa ser traduzida em informação compreensível, contextualizada e pedagogicamente acionável, preservando a possibilidade de interpretação e intervenção do professor.

\section{Tecnologia como Apoio à Decisão Pedagógica}

As tecnologias computacionais podem processar grandes volumes de registros, identificar regularidades e apresentar visualizações que auxiliem o acompanhamento de turmas. Seu valor educacional não decorre apenas da acurácia de um modelo, mas da qualidade dos dados, da interpretação das saídas e da forma como a informação é incorporada à prática docente.

A adoção dessas tecnologias requer que o professor permaneça responsável por interpretar as evidências e decidir intervenções. As previsões precisam explicitar incertezas e limitações, e dados socioeconômicos ou comportamentais não devem ser empregados de modo a naturalizar desigualdades ou reduzir estudantes a categorias fixas. As recomendações também precisam permanecer alinhadas ao currículo e aos objetivos da atividade, enquanto privacidade, consentimento e finalidade de uso dos dados devem ser considerados desde o planejamento.

\section{Técnicas Computacionais na Educação}

\subsection{Aprendizado de Máquina e Inteligência Artificial}

O aprendizado de máquina é uma área da Inteligência Artificial dedicada ao desenvolvimento de modelos capazes de identificar padrões a partir de dados. Em contextos educacionais, os estudos analisados empregaram esses modelos para estimar desempenho ou níveis de proficiência \cite{Multimodels2020_002, Machine2022_010, Math2021_001}, reconhecer padrões de dificuldade, identificar variáveis associadas aos resultados e apoiar recomendações de conteúdos e atividades \cite{Design2025_004}.

Árvores de decisão, florestas aleatórias, máquinas de vetores de suporte (SVM), redes neurais e outros modelos possuem pressupostos e limitações diferentes. A escolha de um algoritmo deve ser comparada a referências simples, considerar o equilíbrio entre as classes e usar métricas adequadas ao problema, em vez de se apoiar apenas na maior acurácia observada.

\subsection{Analítica da Aprendizagem e Mineração de Dados Educacionais}

A analítica da aprendizagem (LA) envolve medir, coletar, analisar e apresentar dados sobre estudantes e seus contextos para compreender e apoiar a aprendizagem. A mineração de dados educacionais (EDM) concentra-se no desenvolvimento e na aplicação de métodos de análise a dados educacionais \cite{Romero2020}.

Essas áreas podem produzir painéis, alertas e modelos de conhecimento, mas a existência de uma correlação não prova uma relação causal. Um padrão de acesso ou tempo de resposta pode refletir diferentes condições, como conhecimento prévio, familiaridade com a interface, acessibilidade ou disponibilidade de conexão.

\subsection{Sistemas Tutores e Aprendizagem Adaptativa}

Sistemas Tutores Inteligentes procuram adaptar a interação ao estado estimado do estudante. A literatura descreve esses sistemas por características como representação do estudante, mecanismos de adaptação, estratégias de orientação e formas de interação \cite{Mousavinasab2021}. Sistemas adaptativos também podem selecionar conteúdos ou ajustar a dificuldade conforme os registros de interação \cite{Xie2019}.

A adaptação é pedagogicamente relevante quando os critérios de decisão são compreensíveis e quando o estudante continua exposto a desafios que promovam autonomia. Uma personalização excessiva ou baseada em dados incompletos pode limitar oportunidades de aprendizagem e reforçar classificações anteriores.

\section{Avaliação Automatizada e Métricas}

Sistemas automatizados podem organizar respostas, calcular indicadores e aplicar regras previamente definidas. Entre as métricas mais comuns estão acurácia, precisão, revocação, medida F1 e área sob a curva ROC \cite{Romero2020}. Essas métricas avaliam o comportamento do modelo e não medem diretamente a qualidade pedagógica da intervenção.

No domínio educacional, a análise também pode considerar acertos e erros por habilidade, tempo e número de tentativas, progressão entre avaliações, participação e conclusão de atividades, além de explicações ou estratégias registradas pelo estudante. A interpretação desses indicadores depende do instrumento, da população e da finalidade. Um modelo com bom desempenho médio pode apresentar erros sistemáticos para determinados grupos; por isso, a avaliação técnica deve incluir análise por classe, validação em dados não utilizados no treinamento e investigação de possíveis vieses.

\section{Requisitos Derivados da Fundamentação}

A fundamentação teórica orientou diretamente a especificação do protótipo. A distinção entre desempenho e aprendizagem levou ao requisito de apresentar indicadores como evidências parciais e contextualizadas; a relação com o currículo sustentou a associação dos resultados a descritores e objetivos educacionais; e a centralidade da mediação docente fundamentou os requisitos de explicabilidade, comunicação de incertezas e revisão humana. Da mesma forma, as discussões sobre dados educacionais justificaram requisitos de privacidade, rastreabilidade e avaliação técnica antes de qualquer alegação de eficácia pedagógica.

Essas relações são consolidadas no Capítulo~\ref{cap:prototipo}, no qual os fundamentos conceituais são convertidos em decisões de projeto sem duplicar a discussão teórica apresentada neste capítulo.

%--------- FIM FUNDAMENTAÇÃO TEÓRICA ------------
